\documentclass[11pt]{article}
\usepackage[margin=1in]{geometry}
\usepackage{amsmath,amsthm,amssymb}
\usepackage[T1]{fontenc}
\usepackage{lmodern}
\newtheorem{theorem}{Theorem}
\newtheorem{lemma}{Lemma}
\newtheorem{corollary}{Corollary}
\newtheorem{remark}{Remark}
\title{D5 230: actual-needed $T0$ core proved uniformly for odd $m$ from the surfaced bundle semantics}
\author{}
\date{2026-03-21}
\begin{document}
\maketitle

\section*{Purpose}
The sharpened background $T0$ from \texttt{d5\_226} still names the full
one-label tiny-repair theorem suggested by artifact \texttt{032}.  But inside the
current front-end manuscript, only a much smaller slice is actually needed:
\begin{itemize}
\item at the mixed source
\[
\sigma_{21}=[4,1,4]/[1,2,2],
\]
remove the one-label tiny-repair branches in
\[
L2,\ L3,\ R2,\ R3;
\]
\item at the canonical bridge seed
\[
\Sigma_{\mathrm{can}}=[2,2,1]/[1,4,4],
\]
remove the one-label tiny-repair branches in
\[
L3,\ R2,\ R3.
\]
\end{itemize}

The goal of this note is to prove exactly this \emph{actual-needed $T0$ core}
for every odd $m>2$ with no remaining background dependency.
The proof uses only the surfaced semantics of
\texttt{torus\_nd\_d5\_endpoint\_latin\_repair.py} and the shared
\texttt{mixed\_008} common module that was already admitted in the
$\sigma_{23}$ and $\sigma_{32}$ closures.

\section*{The repair scope to be killed}
Artifact \texttt{032} allows exactly the following one-label tiny repairs.
\begin{itemize}
\item A \emph{one-gate} repair reroutes one chosen label class to one alternate
absolute direction.
\item A \emph{one-bit} repair reroutes one chosen bit-slice inside one chosen label
class to one alternate absolute direction.
\item A cocycle-defect toggle may additionally act in one of
\[
\{\mathrm{none},\mathrm{left},\mathrm{right},\mathrm{both}\}.
\]
\end{itemize}

Crucially, in the surfaced builder the repair loop only modifies indices whose
current label is exactly the chosen repair label, while the cocycle-defect loop
acts only on the precomputed omitted-layer-$3$ sets.

\section*{A universal surviving $R1$/background collision family}
Throughout we work in the representative color-$0$ chart with
\[
(w_0,s_0)=(0,0).
\]
For the two seed rows relevant here, namely
\[
[2,2,1]/[1,4,4]
\qquad\text{and}\qquad
[4,1,4]/[1,2,2],
\]
the right first-step direction is the same:
\[
\text{right\_word}[0]=1.
\]

\begin{lemma}[universal $R1$/background family]
\label{lem:universal-R1-family}
Fix any modulus $m>2$.  Let
\[
\tau = L/R
\]
be any representative seed in the surfaced builder with
\[
(w_0,s_0)=(0,0),
\qquad
R_1\text{-direction}=1.
\]
Then the color-$0$ row of $\tau$ contains the explicit family
\[
\mathcal F_{R1}(m)
=
\Bigl\{
(x,y,z):
\begin{array}{l}
 x=(x_0,m-1,0,v,u),\\
 y=(x_0,0,0,v,u-1),\\
 z=(x_0,0,0,v,u),\\
 u\in\{1,\dots,m-1\},\\
 x_0+v+u\equiv 2\pmod m
\end{array}
\Bigr\}
\]
of pairwise distinct collision triples with
\[
\mathrm{label}(x)=R1,
\qquad
\mathrm{label}(y)=B,
\qquad
x\mapsto z,
\qquad
y\mapsto z.
\]
In particular, the color-$0$ row is non-Latin.
\end{lemma}

\begin{proof}
Because $u\neq 0$, the source
\[
x=(x_0,m-1,0,v,u)
\]
lies in the layer-$1$ set with
\[
q=m-1,
\qquad
w=0,
\qquad
s=w+u=u\neq 0.
\]
So in the builder notation it is a genuine $R1$ source.  Since the right first
step is direction $1$, the builder sends $x$ to
\[
z=x+e_q=(x_0,0,0,v,u).
\]

Now consider
\[
y=(x_0,0,0,v,u-1).
\]
Its layer is
\[
x_0+0+0+v+(u-1)\equiv 1 \pmod m,
\]
so $y$ is a layer-$1$ source.  But $q(y)=0$, hence $y$ is not in any of the
changed label classes $L1,R1,L2,R2,L3,R3$; therefore its builder label is $B$.

The surfaced common module \texttt{anchor\_by\_feature} assigns anchor $4$ to every
layer-$1$ background feature.  So the baseline row sends $y$ in direction $4$,
that is,
\[
y\mapsto y+e_u=(x_0,0,0,v,u)=z.
\]
Thus $x$ and $y$ are distinct sources with the same image $z$.

Finally, the number of such triples is exactly $m(m-1)$: choose arbitrary
$x_0\in\mathbb Z_m$ and arbitrary nonzero $u\in\mathbb Z_m$, and then $v$ is forced by
\[
x_0+v+u\equiv 2 \pmod m.
\]
So the family is nonempty and explicit for every $m>2$.
\end{proof}

\section*{The actual-needed $T0$ core}
Define the actual-needed repair-label sets by
\[
A_{\mathrm{can}}=\{L3,R2,R3\},
\qquad
A_{21}=\{L2,L3,R2,R3\}.
\]

\begin{theorem}[actual-needed $T0$ core, all odd $m$]
\label{thm:T0-actual-needed-core}
For every odd modulus $m>2$, the following holds.
\begin{enumerate}
\item At the canonical bridge seed
\[
\Sigma_{\mathrm{can}}=[2,2,1]/[1,4,4],
\]
any one-gate or one-bit repair supported on a label in $A_{\mathrm{can}}$, with any
cocycle-defect choice in
\[
\{\mathrm{none},\mathrm{left},\mathrm{right},\mathrm{both}\},
\]
still leaves an explicit collision in color $0$ and is therefore non-Latin.

\item At the mixed source
\[
\sigma_{21}=[4,1,4]/[1,2,2],
\]
any one-gate or one-bit repair supported on a label in $A_{21}$, with any
cocycle-defect choice in the same four-element set, still leaves an explicit
collision in color $0$ and is therefore non-Latin.
\end{enumerate}
\end{theorem}

\begin{proof}
Both seed rows have right first-step direction $1$, so
Lemma~\ref{lem:universal-R1-family} gives the full collision family
$\mathcal F_{R1}(m)$ in color $0$.

It remains to check that the allowed one-label tiny repairs in the statement do
not alter this family.

First, every repair label in $A_{\mathrm{can}}\cup A_{21}$ is different from $R1$.
Hence the repair loop never modifies any $R1$ source $x$ from
$\mathcal F_{R1}(m)$.

Second, the background mate $y$ in Lemma~\ref{lem:universal-R1-family} has label
$B$, so it is never touched by a one-gate or one-bit repair on a non-$B$ label.

Third, the cocycle-defect toggle only acts on the precomputed omitted-layer-$3$
left/right sets.  Both $x$ and $y$ lie in layer $1$, so neither source belongs to
those omitted-layer-$3$ sets.  Therefore the cocycle-defect toggle does not alter
$x\mapsto z$ or $y\mapsto z$.

So for every repair listed in the theorem, the same two sources $x,y$ still map to
$z$.  The color-$0$ row remains non-Latin, hence the repaired candidate is dead.
\end{proof}

\begin{corollary}[front-end consequence]
\label{cor:no-remaining-T0-dependency}
The only $T0$ content actually needed by the current front-end manuscript is now
internalized.
More precisely:
\begin{itemize}
\item at $\sigma_{21}$, every one-label tiny-repair branch in
$L2,L3,R2,R3$ is closed uniformly for odd $m$;
\item at the canonical bridge seed, every one-label tiny-repair branch in
$L3,R2,R3$ is closed uniformly for odd $m$.
\end{itemize}
Therefore the present uses of ``after tiny repair'' inside $T1$ and $T4$ no longer
need an external background theorem.
\end{corollary}

\begin{remark}
This note proves the \emph{actual-needed} $T0$ slice, not the full stronger claim
that every one-label repair across the entire seed list of artifact \texttt{032} has
already been internalized theorem-by-theorem.
For the current manuscript logic, however, the corollary is exactly the dependency
that mattered.
\end{remark}

\section*{Sanity check on the surfaced bundle}
A lightweight checker
\texttt{d5\_230\_T0\_needed\_core\_checker.py} was also written.
On the surfaced bundle it exhaustively tests the actual-needed candidate list for
\[
m=3,5,
\]
and confirms that every tested one-gate / one-bit / cocycle candidate still
preserves the full universal $R1$/background collision family.
The saved summary reports:
\begin{itemize}
\item canonical case: $2016$ checked candidates, all preserving the full family;
\item $\sigma_{21}$ case: $2688$ checked candidates, all preserving the full family.
\end{itemize}
This is only a sanity check; Theorem~\ref{thm:T0-actual-needed-core} itself is the
all-odd proof.

\end{document}
